\section{División del trabajo en robótica entre China y Estados Unidos: hardware, convicción y volante de datos}

Esta sección vuelve a la perspectiva de China y Estados Unidos. Tan Jie (谭杰) considera que el hardware chino avanza muy rápido y que empresas como Yushu (宇树), Zhiyuan (智元) y Xinghaitu (星海图) representan las fortalezas de China en el cuerpo del robot, la cadena de suministro y la manufactura; Estados Unidos/Silicon Valley es más fuerte en la inversión a largo plazo, el cerebro del modelo, la convicción en la investigación y la capacidad de cómputo. En un escenario ideal, China y Estados Unidos deberían colaborar mejor: el desarrollo de la inteligencia en Estados Unidos y la manufactura de hardware en China se complementan.

\begin{figure}[H]
\centering
\includegraphics[width=0.94\textwidth]{china-us-robotics-split.png}
\caption{División del trabajo en robótica entre China y Estados Unidos: Estados Unidos es más fuerte en el paradigma de la inteligencia; China, en hardware y cadena de suministro. Diagrama conceptual de elaboración propia, basado en la conversación de 01:17:35--02:06:16.}
\end{figure}

\begin{knowledgebox}{Lectura de la figura: un hardware fuerte no equivale automáticamente a un volante de datos fuerte}
China es fuerte en hardware, cuerpo del robot, cadena de suministro y manufactura; Estados Unidos/Silicon Valley es fuerte en investigación a largo plazo, el cerebro del modelo, la capacidad de cómputo y la disposición a invertir en visiones de largo plazo. La robótica necesita combinar ambas cosas, pero primero el hardware debe tener un uso real para que el volante de datos empiece a girar.
\end{knowledgebox}

\subsection{La convicción de Silicon Valley y la presión de ciclos cortos en China}

Esta subsección traslada la discusión de la ruta tecnológica al entorno de recursos. La robótica requiere datos e inversión de ingeniería a largo plazo; por ello, la paciencia del capital, la capacidad de hardware y la cultura organizacional de cada región influyen directamente en que una ruta logre concretarse o no.

Tan Jie dice que Silicon Valley está dispuesto a creer en direcciones de largo plazo que parecen ambitious y que no dan resultados a corto plazo, y puede invertir dinero durante diez años; en China se tiende más a la aplicación práctica a corto plazo, la rentabilidad y el crecimiento rápido. Tanto la robótica como los modelos grandes requieren gastar mucho dinero, recolectar muchos datos y realizar experimentos de largo plazo. Si al principio solo se asignan pocos recursos, es muy difícil verificar el scaling.

\begin{warningbox}{La presión de ciclos cortos perjudica la validación de largo plazo}
Muchas rutas de la robótica no pueden validarse con el esquema de «primero dar unos pocos datos para ver el efecto y luego decidir si se invierte». Si el volumen inicial de datos es demasiado pequeño, el modelo ni siquiera puede mostrar las leyes de scaling, y el equipo podría abandonar por error una dirección de largo plazo que era correcta.
\end{warningbox}

\subsection{La ruta de Google, la ruta de Waymo y la ruta de Tesla}

Zhang Xiaojun (张小珺) propone una analogía: la ruta de robótica de Google podría parecerse a la de Waymo, con énfasis en el cerebro y el sistema; las empresas chinas de hardware podrían parecerse a Tesla, con la expectativa de que el despliegue de hardware genere datos. Tan Jie advierte que los autos de Tesla son útiles por sí mismos y por eso su volante de datos puede girar; actualmente, mucho hardware robótico aún no alcanza el umbral de utilidad, por lo que no puede compararse sin más con Tesla.

\begin{knowledgebox}{Condiciones para que funcione el volante de datos}
\begin{tabularx}{\textwidth}{p{0.24\textwidth}p{0.34\textwidth}X}
\toprule
Condición & Significado & Dificultad en robótica \\
\midrule
El producto es útil & Los usuarios están dispuestos a usarlo de verdad, y no solo a ver una demo & Muchos robots aún no alcanzan el umbral de utilidad cotidiana. \\
Escenarios de alta frecuencia & La frecuencia de uso basta para generar una gran cantidad de datos & Los entornos domésticos y abiertos tienen demasiados casos de cola larga. \\
Retroalimentación registrable & Los éxitos, los fracasos y las intervenciones humanas pueden registrarse de forma estructurada & La causa de una falla de manipulación puede estar en la visión, el control de fuerza, la planificación o el hardware. \\
Datos de retorno entrenables & Los datos pueden entrar en el ciclo cerrado de entrenamiento y evaluación & Las diferencias entre cuerpos de robot y la privacidad/seguridad aumentan el costo. \\
\bottomrule
\end{tabularx}
\end{knowledgebox}

\subsection{Resumen de la sección}

La división del trabajo en robótica entre China y Estados Unidos no consiste en que uno reemplace al otro, sino en la complementariedad entre hardware e inteligencia. El verdadero reto es llevar el hardware a escenarios útiles, formar el volante de datos y, con inversión de largo plazo, conectar los modelos, los datos sintéticos y la evaluación en condiciones reales.
